\part{De la scène à une grille de nombres}

\section{Une image n'est pas seulement un tableau RGB}
\subsection{Quatre niveaux de description}
Une photographie existe successivement comme phénomène lumineux,
mesure électronique, représentation numérique et rendu perceptuel.
Le mot « image » désigne souvent l'un de ces niveaux sans le préciser.

\begin{longtable}{p{3cm}p{5.1cm}p{6.5cm}}
\toprule Niveau & Question & Exemple \\ \midrule
Scène physique & Quelle lumière atteint le capteur ? & Un spectre lumineux,
une durée d'exposition et une distribution spatiale. \\
Signal mesuré & Que mesure chaque photosite ? & Une charge électrique,
puis un code numérique filtré par une mosaïque colorée. \\
Image interprétée & Que signifient les composantes ? & RGB linéaire de
caméra, XYZ, sRGB encodé, profondeur, alpha. \\
Fichier & Comment les valeurs sont-elles stockées ? & TIFF avec offsets,
JPEG avec coefficients, PNG avec blocs compressés. \\
Affichage & Quelle lumière l'écran émet-il ? & Conversion vers le profil
écran, fonction de transfert et limites de luminance. \\ \bottomrule
\end{longtable}

Le même tableau $(128,64,32)$ peut désigner une couleur sRGB, des codes
RGB d'un autre espace ou trois mesures sans signification colorimétrique.
Le même fichier peut contenir plusieurs tableaux : miniature, aperçu,
mosaïque RAW, masque, carte de profondeur ou image HDR auxiliaire.

\resultat{L'architecture complète est : des valeurs, une géométrie,
une signification physique ou colorimétrique, une organisation binaire
et des règles d'interprétation. Il manque de l'information si l'un de
ces éléments est absent.}

\subsection{Image matricielle et image vectorielle}
Une image matricielle est définie sur une grille. Une image vectorielle
décrit des objets : chemins, courbes, texte, surfaces. Un SVG peut dire
« tracer un cercle » ; un PNG doit fournir ou reconstruire les valeurs
de ses pixels. La rasterisation transforme le premier en grille à une
résolution choisie. Un PDF peut mêler texte, vecteurs et photographies.
L'agrandissement d'un vecteur recalcule sa couverture ; l'agrandissement
d'une photo interpole des mesures existantes. Une IA peut inventer une
estimation plausible de détails, pas retrouver systématiquement ceux
qui n'ont jamais été mesurés.

\section{Physique de la formation de l'image}
\subsection{Photons, optique et temps d'exposition}
La lumière peut être décrite par un champ électromagnétique ou par des
photons. L'énergie d'un photon de longueur d'onde $\lambda$ est
\begin{equation}
E_\gamma=\frac{hc}{\lambda},
\end{equation}
où $h$ est la constante de Planck et $c$ la vitesse de la lumière.
À environ $550$ nm, cette énergie vaut $3{,}61\times10^{-19}$ joule.
Un pixel ne stocke pourtant pas une liste de photons : le capteur
intègre leur effet pendant l'exposition.

Un modèle simplifié du nombre moyen d'électrons du photosite $i$ est
\begin{equation}
\mu_i=t\int_{A_i}\int E(x,y,\lambda)\,
\eta_i(\lambda)\,\frac{\lambda}{hc}\,d\lambda\,dA.
\end{equation}
$E$ est ici une irradiance spectrale sur le capteur, $A_i$ l'aire du
photosite, $t$ la durée et $\eta_i$ l'efficacité de conversion. Les filtres
colorés et la transmission optique peuvent être inclus dans $E$ ou
dans une réponse effective $\eta_i$. Ce modèle explique trois faits :
une exposition plus longue augmente le signal, les filtres changent
la réponse spectrale et les photosites ne mesurent pas tous la même
bande de lumière.

Pour une surface diffuse, un modèle de scène souvent utile est
$E(\lambda)\propto S(\lambda)\rho(\lambda)$ : spectre de l'éclairage
fois réflectance. Il n'est pas complet pour les surfaces spéculaires,
translucides ou fluorescentes. Une couleur de peinture publiée sur un
site marchand ne donne ni son spectre de réflectance ni l'éclairage réel.

\subsection{Flou optique et fonction d'étalement}
Même avant la conversion numérique, un point de la scène devient une
tache. Dans un modèle local linéaire invariant, l'image idéale $f$ est
convoluée par une fonction $h$, dite PSF :
\begin{equation}
g(x,y)=(h*f)(x,y)=\iint h(u,v)f(x-u,y-v)\,du\,dv.
\end{equation}
Défocalisation, diffraction et mouvement participent au flou. Dans
le domaine fréquentiel, la convolution devient un produit ; la MTF
décrit la transmission des contrastes suivant la fréquence. Si une
fréquence est supprimée et noyée dans le bruit, une accentuation ne
la recrée pas avec certitude. Elle augmente surtout le contraste de
ce qui subsiste.

\subsection{Exposition et ISO : deux actions différentes}
À scène et géométrie constantes, doubler $t$ double approximativement
les photons collectés tant qu'il n'y a pas saturation. Une variation
de $e$ EV représente un facteur $2^e$. Pour un objectif idéal et un
champ donné, l'éclairement varie approximativement comme $1/N^2$, où
$N$ est le nombre d'ouverture. Passer de $f/4$ à $f/2{,}8$ représente
environ un doublement de l'éclairement.

L'ISO règle une chaîne de gain et de présentation de la sensibilité,
selon l'appareil ; ce n'est pas une commande qui multiplie les photons.
Le gain peut être analogique, numérique ou associé à des modes de
lecture différents. « ISO plus élevé » ne signifie donc pas
automatiquement « moins de photons » : ce sont les choix d'exposition
associés qui le déterminent. Pour comparer le bruit, documenter le
temps, l'ouverture, l'éclairage, l'ISO et le mode RAW.

\section{Capteur, bruit et conversion analogique-numérique}
\subsection{Modèle simplifié d'un code RAW}
Un modèle pédagogique est
\begin{equation}
D_i=\operatorname{clip}\!\left(
\operatorname{round}(B_i+K_i(N_i+n_{r,i})),0,2^b-1\right),
\end{equation}
avec $B_i$ un offset noir en codes, $K_i$ un gain en codes par électron,
$N_i$ les électrons photo-générés, $n_{r,i}$ un bruit de lecture et $b$
la profondeur de conversion. Des traitements internes et des
non-linéarités peuvent compléter ce modèle. La norme de mesure
EMVA~1288 distingue notamment caractérisation linéaire et générale
des capteurs~\cite{emva}. Ce cours utilise le modèle simple pour
raisonner ; il ne certifie aucun boîtier.

\subsection{Bruit de photons et bruit de lecture}
Pour un comptage de photons de type Poisson,
\begin{equation}
N\sim\mathcal P(\mu),\qquad \mathbb E[N]=\mu,
\qquad \operatorname{Var}(N)=\mu.
\end{equation}
L'écart-type vaut $\sqrt\mu$ : recevoir davantage de lumière augmente
le bruit absolu, mais améliore le rapport signal/bruit relatif.
Si le bruit de lecture indépendant a une variance $\sigma_r^2$,
\begin{equation}
\operatorname{SNR}\simeq\frac{\mu}{\sqrt{\mu+\sigma_r^2}}.
\end{equation}
Avec $\mu=100$ électrons et $\sigma_r=3$, on obtient environ $9{,}58$.
Avec $\mu=400$, on obtient $19{,}78$. Multiplier ensuite le fichier
sombre par quatre amplifie aussi son bruit et ne procure pas les
400 électrons d'une vraie exposition plus longue.

Le bruit thermique dépend du temps et de la température ; les
inhomogénéités de gain ou d'offset produisent du bruit spatial fixe.
Un pixel chaud, une ligne de lecture et un bruit aléatoire ne se
traitent pas forcément avec le même algorithme.

\subsection{Niveau noir, niveau blanc et saturation}
Le code zéro n'est pas nécessairement le noir du capteur. Un offset
permet notamment de représenter les fluctuations autour du noir
sans les couper trop tôt. Pour une normalisation simple :
\begin{equation}\label{eq:normalisation}
u=\frac{D-B}{W-B},\qquad W>B.
\end{equation}
Si $D=6144$, $B=2048$ et $W=16383$, alors $u\approx0{,}28574$.
Changer arbitrairement $W$ modifie l'échelle de tous les pixels.
Changer $B$ affecte proportionnellement davantage les ombres.

Le blanc opérationnel peut être inférieur au maximum théorique
$2^b-1$. Il dépend parfois du canal, de l'ISO ou du mode. Les marges
optical-black sont des zones masquées destinées à mesurer le noir,
mais toutes les bordures du tableau ne sont pas nécessairement masquées.
Une zone active prise pour du noir produit une correction erronée.

Au-dessus de la capacité de saturation, plusieurs éclairements
distincts donnent le même code. Aucune transformation réversible
du fichier ne peut les distinguer. Une reconstruction de hautes
lumières peut exploiter des canaux non saturés ou une hypothèse de
couleur ; elle n'est pas une récupération certaine de la mesure perdue.

\section{Échantillonnage spatial et résolution}
\subsection{Du champ continu à la grille}
Une image monochrome discrète est une fonction
$I:\{0,\ldots,W-1\}\times\{0,\ldots,H-1\}\to\mathbb R$.
Le pixel représente une mesure spatialement intégrée, pas un point
mathématique infiniment petit. Une grille $6000\times4000$ contient
24 millions de positions. Cela ne garantit ni 24 millions de détails
indépendants ni une résolution optique correspondante.

Une fréquence spatiale trop élevée peut se replier en fréquence
plus basse : c'est l'aliasing. En une dimension, une sinusoïde à
$0{,}7$ cycle par échantillon peut être confondue, sur la grille,
avec une fréquence de $0{,}3$. La limite de Nyquist est $0{,}5$ cycle
par échantillon pour un signal correctement limité en fréquence.
En photographie, les structures fines, tissus et grilles peuvent
produire du moiré ; le CFA ajoute sa propre structure d'échantillonnage.

\subsection{Redimensionner n'est pas seulement sélectionner des pixels}
Réduire une image par deux doit généralement être précédé d'un
passe-bas. Conserver un pixel sur deux sans filtrer peut transformer
une texture fine en motif faux. Agrandir reconstruit une fonction
interpolée : plus proche voisin, bilinéaire, bicubique ou autre.
Le choix change netteté, halos et aliasing ; il n'ajoute pas de
mesures physiques nouvelles.

Un filtre bilinéaire en deux dimensions combine quatre voisins :
\begin{equation}
I(x,y)=(1-a)(1-b)I_{00}+a(1-b)I_{10}
+(1-a)bI_{01}+abI_{11},
\end{equation}
où $a,b\in[0,1]$ sont les fractions de position. Si les valeurs sont
des couleurs encodées non linéaires, cette moyenne n'est pas la
même que la moyenne des lumières ; nous y reviendrons.

\subsection{DPI, PPI, dimensions et taille physique}
Une largeur de 6000 pixels imprimée à 300 pixels par pouce donne
$6000/300=20$ pouces, soit $50{,}8$ cm. Changer seulement la métadonnée
de résolution à 150 PPI ne crée aucun pixel : la dimension physique
demandée devient 101,6 cm. Le DPI d'une imprimante décrit ses points
matériels et n'est pas nécessairement égal au nombre de pixels d'image.

\section{Mosaïque colorée, canaux et profondeur}
\subsection{Photosite, CFA et pixel RGB}
Dans un capteur Bayer courant, chaque photosite mesure une réponse
spectrale principalement rouge, verte ou bleue. Une cellule RGGB est
\[
\begin{matrix}R&G\\G&B\end{matrix}.
\]
Ce sont des filtres larges et chevauchants, pas des monochromateurs
parfaits. Les deux verts occupent des positions différentes ; on
peut devoir préserver quatre classes RGGB dans la calibration.
Le dématriçage estime les composantes absentes à chaque position.
Le RAW mosaïqué de 24 millions de photosites contient ainsi environ
24 millions de mesures, et non 72 millions de mesures RGB indépendantes.

Des capteurs monochromes, des CFA différents, des capteurs à couches
et des modes multivues existent. Il ne faut jamais déduire le motif
de la seule extension du fichier. Le décalage du recadrage peut aussi
changer la phase du motif : enlever une colonne fait commencer une
ligne RG par GR.

\subsection{Profondeur par échantillon et profondeur par pixel}
Un entier non signé sur $b$ bits possède $2^b$ codes, de 0 à $2^b-1$.
Une image RGB 8 bits par composante utilise 24 bits par pixel avant
compression ; RGBA utilise 32 bits. Un RAW 14 bits mosaïqué n'est
donc pas « moins riche » qu'un JPEG 24 bits au sens où ces nombres
seraient directement comparables : ils décrivent des objets différents.

Une quantification uniforme normalisée est
\begin{equation}
Q_b(x)=\frac{\operatorname{round}((2^b-1)x)}{2^b-1},\quad x\in[0,1].
\end{equation}
Le pas est $\Delta=1/(2^b-1)$ et l'erreur d'arrondi est bornée par
$\Delta/2$. Sous des hypothèses statistiques supplémentaires,
sa variance est approximativement $\Delta^2/12$ ; cette erreur n'est
pas systématiquement un bruit aléatoire indépendant. Sur un dégradé
très lisse, elle peut produire des bandes. Le dithering ajoute un
signal contrôlé avant quantification pour réduire ces structures,
au prix d'un bruit visible ou mesurable.

\subsection{Profondeur nominale, utile et stockage}
Un signal 12 bits peut être rangé dans un entier 16 bits. Il faut
savoir s'il occupe les bits faibles, les bits forts ou une
représentation compactée. La profondeur de stockage ne démontre
pas le nombre de bits utiles du capteur. De même, convertir un
JPEG 8 bits en TIFF 16 bits ne recrée pas les valeurs intermédiaires
perdues : cela donne surtout davantage de marge pour les calculs suivants.

\resultat{Résolution spatiale, nombre de composantes, précision numérique,
dynamique physique, gamut et compression sont des dimensions distinctes
de la qualité. Aucun nombre unique ne les résume.}
